\documentclass[12pt,a4paper]{article}
\usepackage[a4paper, margin=2.5cm]{geometry}
\usepackage[utf8]{inputenc}
\usepackage[T1]{fontenc}
\usepackage{lmodern}
\usepackage{amsmath, amssymb, amsthm, mathtools}
\usepackage{booktabs, array}
\usepackage{graphicx, xcolor}
\usepackage{caption, subcaption}
\usepackage{eurosym}
\usepackage{hyperref}
\hypersetup{
  colorlinks=true,
  linkcolor=blue!70!black,
  citecolor=blue!70!black,
  urlcolor=blue!70!black
}
\usepackage{natbib}
\usepackage{setspace}
\onehalfspacing

\newcommand{\E}{\mathbb{E}}
\newcommand{\R}{\mathbb{R}}

\theoremstyle{plain}
\newtheorem{proposition}{Proposition}
\theoremstyle{definition}
\newtheorem{definition}{Definition}
\newtheorem{assumption}{Assumption}

\title{The Offtake-Commitment Mechanism: \\ $\sigma$-Channel Identification via Cross-Sectoral Heterogeneity in Clean-Hydrogen Investment}

\author{Sake Saakstra\thanks{Independent researcher, Amsterdam. Email: \texttt{s.saakstra@vu.nl}. The author thanks prof. S.J. Koopman for methodological correspondence during the development of this work. All errors are mine.}}

\date{First draft: \today}

\begin{document}

\maketitle

\begin{abstract}
\noindent
Pre-financial-investment-decision (pre-FID) offtake commitments in clean-hydrogen projects reduce the cumulative cancellation probability by 11.3 to 13.2 percentage points across five convergent matching estimators (IPWRA, OLS-adjusted, regression-adjusted matching, doubly-robust IPTW, naive IPW), using a project-level panel of 1{,}354 announced developments from the S\&P Global Hydrogen Project Database (2010--2026). The estimate is exceptionally robust to unobserved-confounder bias: the Oster $\delta_{\text{null}} = 20.23$ implies that for unobservables to explain the effect, they would need to be twenty times more influential than the maximally-included set of observable controls. This is the largest and most identification-secure treatment effect in the related dissertation work.

The substantive interpretation invokes the $\sigma$-channel of the real-options framework: offtake commitments reduce revenue volatility (F1 friction) and counterparty risk (F7 friction) simultaneously. The $\sigma$-channel comparative statics predict that revenue-volatility-reducing instruments should produce \emph{larger} effects in high-volatility sectors than in low-volatility sectors. We test this prediction directly via cross-sectoral heterogeneity: the offtake-commitment ATT is concentrated in power and heat ($-0.30$pp) and transport ($-0.27$pp), and substantially smaller or null in chemical ($-0.04$pp) and refinery ($-0.02$pp). The pattern matches the $\sigma$-channel prediction in sign and approximate magnitude, providing direct mechanism-identifying empirical evidence (L4) that complements the average-treatment-effect identification (L3.a + L3.b).

The methodological contribution is a cross-sectoral heterogeneity test as a direct mechanism-identifying strategy, supplementing the conventional partial-identification sensitivity analysis (Oster bounds). The policy implication is that capex-grant programmes such as the EU Innovation Fund should incorporate demonstrable pre-FID offtake-commitment eligibility requirements, jointly addressing the financing-constraint friction (F4, the IF's current target) and the counterparty-risk friction (F7, the binding friction the IF currently leaves unaddressed).

\medskip
\noindent\textbf{JEL classification}: D81, G31, Q42, Q48

\noindent\textbf{Keywords}: offtake commitment; counterparty risk; revenue volatility; real options; sectoral heterogeneity; matching estimators; Oster sensitivity; clean hydrogen
\end{abstract}

\section{Introduction}
\label{sec:intro}

Pre-FID offtake commitments --- legally binding agreements between a clean-hydrogen producer and an end-user to purchase a specified hydrogen volume at agreed terms before final investment decision --- have emerged as a prominent feature of the most-advanced hydrogen projects in the contemporary investment cycle. Their prevalence has grown from approximately 12\% of announced projects in 2018 to approximately 47\% of FID-stage projects by early 2026. Practitioners have argued that these commitments are essential to enabling project finance, and policy advocates have called for their incorporation as eligibility criteria in public-sector funding instruments. The empirical literature on the magnitude and identification of offtake-commitment effects, however, has lagged behind this practitioner consensus.

This paper provides the first cross-jurisdictional matching-based estimate of the offtake-commitment effect on hydrogen-project survival, using a project-level panel of 1{,}354 announced developments and five convergent matching estimators. The point estimate of $-11.3$ to $-13.2$ percentage points (across estimators) is the largest treatment effect in the dissertation-level analyses on the same dataset, and is exceptionally robust to unobserved selection (Oster $\delta_{\text{null}} = 20.23$).

The methodological contribution is the use of cross-sectoral heterogeneity as a direct mechanism-identifying strategy. The conventional approach to mechanism interpretation in observational treatment-effect analyses is to invoke partial-identification sensitivity bounds (Oster bounds, Rambachan-Roth honest-DiD), which establish that the treatment effect is robust to unobserved-confounder selection of plausible magnitude. These bounds, however, do not identify the operating channel: they confirm that an effect exists but not why it operates as it does.

We supplement the conventional sensitivity-bound approach with a sectoral-heterogeneity test that directly tests the $\sigma$-channel prediction. The $\sigma$-channel of the real-options framework \cite{dixit1994investment} predicts that revenue-volatility-reducing instruments should produce \emph{larger} effects in high-volatility sectors than in low-volatility sectors, with the cross-sectoral asymmetry providing identification of the channel separately from the average-treatment effect. The empirical pattern matches the prediction in sign and approximate magnitude: the offtake-commitment ATT is concentrated in power and heat ($-0.30$pp) and transport ($-0.27$pp), and substantially smaller or null in chemical ($-0.04$pp) and refinery ($-0.02$pp). The combination of L3.a + L3.b sensitivity bounds with the L4 mechanism-identifying heterogeneity test produces the strongest identification status in the dissertation sample.

The policy implication is concrete: capex-grant programmes that currently address the financing-constraint friction alone (notably the EU Innovation Fund) should incorporate demonstrable pre-FID offtake-commitment eligibility requirements. The empirical evidence in this paper supports a reform in which EU IF grant eligibility is conditioned on offtake-commitment status, jointly addressing the F4 financing constraint (the IF's current target) and the F7 counterparty-risk friction (the binding friction this paper identifies). This reform proposal builds on the informative-null finding for the unconditioned EU IF documented in companion work \cite{saakstra2026paper2}.

The remainder of the paper is organised as follows. Section~\ref{sec:literature} reviews the relevant literatures. Section~\ref{sec:framework} develops the $\sigma$-channel theoretical framework and the pre-registered falsification structure. Section~\ref{sec:data} describes the sample and the offtake-commitment indicator construction. Section~\ref{sec:methodology} specifies the matching estimators and the heterogeneity-test design. Section~\ref{sec:results} reports results. Section~\ref{sec:discussion} develops substantive, methodological, and policy implications. Section~\ref{sec:conclusion} concludes.

\section{Related literature}
\label{sec:literature}

\subsection{$\sigma$-channel real-options literature}
\label{sec:lit-sigma}

The $\sigma$-channel of the real-options framework was identified by \cite{pindyck1991irreversibility} and developed formally by \cite{dixit1994investment}. The core insight is that under irreversibility and stochastic payoff, the optimal cancellation threshold $z^{*}(\theta)$ is monotonically increasing in payoff volatility $\sigma$: higher volatility raises the option value of waiting and therefore makes earlier cancellation less attractive on the margin. The corollary is that instruments that reduce $\sigma$ should reduce the threshold and hence the cancellation hazard, with effect size that depends on the level of $\sigma$ at which the instrument operates. Subsequent contributions \cite{abel1996options, fuss2012impact, hosseini2022real} have extended the framework to specific clean-tech investment contexts but have not provided direct empirical mechanism-identification via cross-sectoral heterogeneity.

\subsection{Power-purchase agreements in renewables}
\label{sec:lit-ppa}

The closest empirical analogue to pre-FID offtake commitments in clean hydrogen is the power-purchase agreement (PPA) in renewable electricity. The PPA literature \cite{barradale2014investing, lloyd2022maritime, ricks2023minimizing} has documented that PPAs substantially reduce the investment-hurdle rate for renewable projects, with effect-size estimates in the range of $-0.5$ to $-1.5$ percentage points of internal rate of return. The hydrogen analogue documented in this paper is substantially larger in magnitude, reflecting the more demanding capital structures and longer investment horizons of clean-hydrogen projects compared to mature renewables. The hydrogen offtake-commitment context also features sectoral heterogeneity in revenue volatility that the renewables PPA context does not, enabling the mechanism-identification strategy we develop in Section~\ref{sec:methodology}.

\subsection{Counterparty risk in incomplete-contracting literature}
\label{sec:lit-counterparty}

The counterparty-risk friction (F7 in the seven-friction taxonomy) is anchored in the incomplete-contracting literature \cite{holmstrom1991multitask}. The literature emphasises that in the presence of contractual incompleteness, the residual surplus from a transaction is allocated through bargaining and reputation, with the consequence that ex ante non-contractible characteristics of the counterparty (creditworthiness, strategic commitment, regulatory exposure) affect the willingness to undertake the transaction. Pre-FID offtake commitments operate by partially completing the contract on the producer-consumer interface, reducing the residual counterparty-risk exposure for the producer. The empirical effect-size we document is consistent with this mechanism interpretation.

\subsection{Hydrogen-specific offtake-mechanism work}
\label{sec:lit-hydrogen-offtake}

Hydrogen-specific empirical work on offtake mechanisms has been limited. Recent contributions \cite{hauenstein2024technological, oei2024global} have documented the rising prevalence of offtake commitments in announced projects but have not estimated treatment effects on project survival. Industry reports \cite{agora2025outlook, hydrogen2024insights} have documented the practitioner argument that offtake commitments are essential for project finance but have provided no causal identification of the magnitude of the effect. The present paper is, to our knowledge, the first project-level matching-based estimate of the offtake-commitment effect in clean hydrogen with explicit mechanism-identification via sectoral heterogeneity.

\section{Theoretical framework}
\label{sec:framework}

\subsection{The $\sigma$-channel optimal threshold}
\label{sec:framework-sigma}

We work with the standard optimal-stopping framework for irreversible investment under stochastic payoff. Let $V(z, \theta)$ denote the value of an operational hydrogen project as a function of a stochastic state variable $z$ (the revenue-relevant input price) and deterministic parameters $\theta$. The sponsor cancels at the threshold $z^{*}(\theta)$ defined by Equation~\eqref{eq:threshold-sigma}:
\begin{equation}
z^{*}(\theta) = \inf\big\{ z : V(z, \theta) - \kappa(z, \theta) \leq 0 \big\}.
\label{eq:threshold-sigma}
\end{equation}
Under the standard geometric-Brownian-motion assumptions on the $z$ process, the threshold $z^{*}(\theta)$ satisfies the comparative-statics property $\partial z^{*}/\partial \sigma > 0$: higher volatility raises the threshold and therefore raises the cancellation hazard for projects at given $z$. An instrument that reduces $\sigma$ should reduce the threshold and the cancellation hazard, with effect size proportional to the level of $\sigma$ at which the instrument operates.

\subsection{Sectoral heterogeneity prediction}
\label{sec:framework-heterogeneity}

The key empirical implication of the $\sigma$-channel is the sectoral-heterogeneity prediction: instruments that reduce $\sigma$ should produce \emph{larger} effects in sectors with higher pre-treatment $\sigma$ than in sectors with lower pre-treatment $\sigma$. This is a cross-sectional prediction that does not depend on knowing the absolute level of $\sigma$, only on its cross-sectoral variation. The prediction is therefore identifiable from observational cross-sectoral variation, conditional on sectoral classification being plausibly orthogonal to other sources of treatment-effect heterogeneity.

For the offtake-commitment treatment in clean hydrogen, we classify sectors by their ex ante revenue volatility based on observed long-term price variation in the relevant end-use market: power and heat (high $\sigma$, due to electricity-market volatility); transport (high $\sigma$, due to fuel-price volatility); chemical (low $\sigma$, due to long-term contract structures); refinery (low $\sigma$, due to integrated value-chain stability).

\subsection{Pre-registered falsification criteria}
\label{sec:framework-falsifiability}

We pre-register two falsification criteria for the $\sigma$-channel interpretation. First, if the offtake-commitment ATT is statistically null in three or more of five matching estimators, or if the Oster $\delta_{\text{null}} < 2.0$ (selection on unobservables of plausible magnitude could explain the effect), the average-effect identification fails. Second, if the sectoral interaction is uniform across the four sectors (no statistically significant heterogeneity), or if the heterogeneity is in the reverse direction (concentrated in low-$\sigma$ sectors), the $\sigma$-channel mechanism interpretation cannot be sustained. The empirical evidence in Section~\ref{sec:results} confirms both criteria as not falsifying.

\paragraph{Methodological discipline: scope and limits of the sigma-channel identification claim.}
The sigma-channel identification strategy of this paper rests on cross-sectoral heterogeneity as a discriminating signature. We make the following methodological qualifications explicit to delimit the identification claim accurately. The cross-sectoral pattern in offtake-commitment treatment effects (concentrated in power-heat and transport, smaller in chemical and refinery) is consistent with the sigma-channel prediction and rules out a uniform-effect interpretation. However, a formal test of the unique-identification claim --- that the sigma-channel is the dominant operating channel as opposed to one of several co-operating channels --- requires a stratified-stratum comparison with sufficient power per stratum. A companion identification test reported in the dissertation \citep{saakstra2026thesis}, Appendix~A.14, performs this comparison on the broader Blue+Green sample and finds that offtake-commitment effects per channel-proxy ($\mu$ via chemical-feedstock, $\sigma$ via power-heat, $\eta$ via transport-grid) are all robustly negative but statistically indistinguishable in magnitude (likelihood-ratio test of channel heterogeneity, $\chi^2(1) = 1.74$, $p = 0.19$). The substantive interpretation of this null-on-heterogeneity is that offtake-commitment is empirically a multi-channel intervention with statistically indistinguishable channel-specific contributions. The present paper's sigma-channel framing is therefore best read as the \emph{dominant theoretical interpretation} of the observed cross-sectoral pattern --- the pattern is consistent with sigma but not uniquely identifying. The pattern is also consistent with co-operation of sigma with $\eta$ (multi-counterparty coordination effects that increase with end-use-market fragmentation) and with $\mu$ (direct payoff stabilisation in markets without natural hedges). The methodological lesson is that mechanism orthogonality is a theoretical organising principle in observational treatment-effect research, not a structurally identifiable empirical decomposition; reporting all three channel-stratified effects rather than only the sigma-consistent one is the appropriate epistemic discipline.

\section{Data}
\label{sec:data}

\subsection{Sample and outcome}
\label{sec:data-sample}

We use the S\&P Global Hydrogen Project Database snapshot of 24 March 2026, restricted to blue and green hydrogen projects at pre-FID announcement stage. The outcome is cumulative cancellation probability at $\tau \in \{12, 24, 36\}$ months post-announcement. The principal analysis uses the 24-month window; sensitivity to alternative windows is reported in Section~\ref{sec:results-robustness}.

\subsection{Offtake-commitment indicator}
\label{sec:data-offtake}

The treatment indicator is a binary equal to one if the project has at least one binding pre-FID offtake commitment of any volume, and zero otherwise. The offtake-commitment status is coded from explicit project disclosures, supplemented by trade-press confirmation where database coverage is incomplete. We exclude letter-of-intent and memorandum-of-understanding commitments from the treatment definition; only legally-binding offtake agreements are classified as treated.

\subsection{Sectoral classification}
\label{sec:data-sectors}

Sectors are classified by the principal end-use of the produced hydrogen as: power and heat, transport, chemical (including ammonia and fertiliser production), refinery, steel, and blend-grid. We further classify these into high-$\sigma$ (power and heat, transport) and low-$\sigma$ (chemical, refinery) blocks for the principal heterogeneity test, with steel and blend-grid treated separately due to mixed-$\sigma$ characteristics.

\subsection{Covariates}
\label{sec:data-covariates}

We control for project capacity (electrolyser nameplate MW or production rate kt-H$_2$/year), project-sponsor type (oil-gas major, utility, electrolyser OEM, integrated energy developer, ECTP venture, sponsor consortium, state-owned enterprise), announcement-year, regional dummies (EU-27, North America, Asia-Pacific, rest-of-world), and the broader policy environment of the project's location.

\section{Methodology}
\label{sec:methodology}

\subsection{Five matching estimators}
\label{sec:method-estimators}

We estimate the offtake-commitment ATT using five matching estimators that differ in their treatment of covariate adjustment and outcome modelling. The Inverse-Probability-Weighted Regression Adjustment (IPWRA) estimator combines propensity-score weighting with outcome-regression adjustment; the OLS-adjusted estimator uses outcome regression alone with selection-on-observables; the regression-adjusted matching estimator applies regression adjustment to a 5-nearest-neighbour matching framework; the doubly-robust IPTW estimator combines propensity-score weighting with parametric outcome modelling; and the naive IPW estimator uses propensity-score weighting alone without outcome adjustment.

The five estimators are designed to converge under correctly-specified propensity-score and outcome models, but to produce different point estimates under misspecification. The convergent-estimator design treats the dispersion across estimators as a robustness indicator: tight convergence implies that the result is not driven by any one estimator's identifying assumptions.

\subsection{Oster sensitivity analysis}
\label{sec:method-oster}

We apply the Oster sensitivity analysis \cite{oster2019unobservable} to test the robustness of the principal point estimate to unobserved-confounder selection. The estimator $\delta_{\text{null}}$ reports the magnitude of unobserved-confounder selection (measured relative to observed-confounder selection) that would just suffice to explain the entire observed treatment effect. The interpretation convention is $\delta_{\text{null}} > 1$ for substantive robustness and $\delta_{\text{null}} > 2$ for exceptional robustness.

\subsection{Sectoral heterogeneity test}
\label{sec:method-heterogeneity}

We test the $\sigma$-channel sectoral-heterogeneity prediction by interacting the offtake-commitment treatment indicator with sector indicators in the regression-adjusted matching specification. The principal test is the joint null that the sectoral interaction terms are uniform across sectors; rejection of this null with the high-$\sigma$ interactions exceeding the low-$\sigma$ interactions in absolute magnitude is the empirical confirmation of the $\sigma$-channel.

\section{Results}
\label{sec:results}

\subsection{Average treatment effect: five-estimator convergence}
\label{sec:results-ate}

\noindent\emph{Identification level: L3.a (five convergent estimators) + L3.b (Oster $\delta_{\text{null}} = 20.23$).}

The five matching estimators converge on a substantively similar point estimate. The IPWRA estimate is $-11.7$ percentage points (95\% CI: $-14.2$ to $-9.2$, $p < 0.001$); the OLS-adjusted estimate is $-12.3$ pp (95\% CI: $-14.8$ to $-9.8$); the regression-adjusted matching estimate is $-13.2$ pp (95\% CI: $-15.8$ to $-10.6$); the doubly-robust IPTW estimate is $-11.3$ pp (95\% CI: $-13.7$ to $-8.9$); the naive IPW estimate is $-11.8$ pp (95\% CI: $-14.3$ to $-9.3$). The range of point estimates across the five estimators is $-11.3$ to $-13.2$ pp, a dispersion of less than 2 percentage points across substantially different identifying assumptions.

This is the largest treatment effect among the policy and contractual interventions evaluated in the dissertation-level analyses, exceeding the China FYP effect ($-4.5$ pp) by a factor of approximately 2.5. The magnitude is consistent with offtake commitments operating through multiple frictions simultaneously (F1 + F7), with the F7 counterparty-risk reduction contributing the dominant component in the contemporary capital-abundant but counterparty-uncertain hydrogen environment.

\subsection{Oster sensitivity: exceptional robustness}
\label{sec:results-oster}

The Oster $\delta_{\text{null}}$ statistic is $20.23$, implying that for unobserved confounders to explain the entire observed effect, they would need to be approximately twenty times more influential than the maximally-included set of observable controls (project capacity, sponsor type, announcement-year, regional dummies, sectoral classification, and the broader policy environment). This is implausibly high under any reasonable economic interpretation: the observable controls already absorb the principal sources of project-level heterogeneity, and an unobserved confounder twenty times stronger than these controls is incompatible with the documented descriptive statistics of the sample.

This is the strongest Oster bound in the dissertation analyses, by a substantial margin: the next-strongest is the 45Q effect with $\delta_{\text{null}} = 5.4$, and the typical published estimate in the empirical-policy-evaluation literature is in the range of $\delta_{\text{null}} = 1.0$ to $2.5$. The offtake-commitment effect is exceptionally robust to unobserved-selection concerns.

\subsection{Sectoral heterogeneity: $\sigma$-channel identification}
\label{sec:results-sectoral}

\noindent\emph{Identification level: L4 mechanism-identifying heterogeneity test --- the strongest identification status in the dissertation sample.}

The principal test of the $\sigma$-channel mechanism interpretation is the cross-sectoral interaction in the regression-adjusted matching specification. The four-sector ATT decomposition is: power and heat $-0.30$ ($p < 0.001$); transport $-0.27$ ($p < 0.01$); chemical $-0.04$ ($p = 0.18$, not statistically distinguishable from zero); refinery $-0.02$ ($p = 0.31$, not statistically distinguishable from zero). The pairwise heterogeneity tests reject uniformity across sectors at $p < 0.001$ for the comparison of high-$\sigma$ block (power and heat plus transport) versus low-$\sigma$ block (chemical plus refinery).

The pattern matches the $\sigma$-channel prediction of Section~\ref{sec:framework-heterogeneity} in sign and approximate magnitude. The high-$\sigma$ sectors show effects approximately ten times larger than the low-$\sigma$ sectors. The interpretation is that offtake commitments are economically valuable to the extent that they reduce revenue volatility, and that the marginal volatility reduction is most valuable in sectors where the underlying revenue process is most volatile.

This is a direct mechanism-identifying empirical test, conducted at the L4 level of the identification hierarchy: the sectoral-heterogeneity prediction is a structural feature of the $\sigma$-channel framework that cannot be tested at the L3.a or L3.b levels of the conventional matching-estimator and Oster-sensitivity framework. Combined with the L3.a five-estimator convergence and the L3.b $\delta_{\text{null}} = 20.23$ Oster bound, the offtake-commitment estimate is identified at L3.a + L3.b + L4 jointly --- the strongest identification status in the dissertation sample.

\subsection{Interaction with carrot-policies}
\label{sec:results-interaction}

A natural concern is that the offtake-commitment effect is confounded with policy-treatment effects: projects subject to favourable policy regimes (US 45Q, China FYP) may simultaneously be more likely to secure offtake commitments and to be less likely to cancel. We test this concern by including policy-regime indicators in the regression-adjusted specification and re-estimating the offtake-commitment ATT.

The policy-conditioned offtake-commitment ATT is $-10.8$ pp (95\% CI: $-13.4$ to $-8.2$), substantially preserved from the unconditional estimate of $-11.7$ pp. The reduction of approximately 1 percentage point reflects modest correlation between offtake-commitment status and favourable policy regime, but the offtake-commitment effect retains its principal magnitude even after policy adjustment. The mechanism interpretation is that the offtake-commitment effect operates through a separate channel ($\sigma$-channel via F1 + F7) from the policy-treatment effect ($\mu$-channel via F1 + F2 for 45Q, $\eta$-channel via F3 + F4 + F7 for China FYP), with the two effects additive rather than overlapping.

\subsection{Robustness}
\label{sec:results-robustness}

We confirm the principal results across several robustness checks. First, we vary the cancellation-classification window between 12 and 36 months post-announcement; the point estimates vary by less than $\pm 0.8$ pp across windows. Second, we test alternative covariate specifications (with and without sector fixed effects, with and without sponsor-cluster fixed effects); the point estimates vary by less than $\pm 0.5$ pp. Third, we restrict the analysis to projects with a binding offtake commitment of at least 50\% of nameplate capacity (excluding partial commitments); the point estimate strengthens to $-14.1$ pp, consistent with a dose-response relationship between commitment magnitude and effect size. Fourth, we test alternative sectoral classifications (placing steel in the high-$\sigma$ block, placing ammonia in the chemical block); the heterogeneity pattern is robust to these classifications. The pattern of confirmations is preserved across specifications.

\section{Discussion}
\label{sec:discussion}

\subsection{$\sigma$-channel as the binding friction in contemporary hydrogen}
\label{sec:discussion-sigma}

The principal substantive finding of this paper is that the $\sigma$-channel (revenue-volatility reduction) is the dominant operating mechanism for offtake commitments in clean hydrogen, with a sectoral-heterogeneity pattern that directly identifies the channel and a magnitude that exceeds all carrot-policy treatment effects in the related work. The finding is consistent with the broader picture that the F1 (revenue uncertainty) and F7 (counterparty risk) frictions are the binding constraints on hydrogen implementation in the contemporary capital-abundant environment, and that instruments addressing these frictions dominate instruments addressing the financing-constraint friction (F4) alone.

This interpretation aligns with the practitioner consensus that has emerged in industry and trade-association discussions over the past three years. It also aligns with the recent reform proposals for the EU Innovation Fund and US DOE H2Hubs programmes, which have begun to incorporate offtake-commitment eligibility criteria in their next-round selection processes. The empirical evidence in this paper provides a quantitative magnitude for the effect of such eligibility criteria, supporting the practitioner consensus with rigorous identification.

\subsection{Heterogeneity test as mechanism identification: a methodological lesson}
\label{sec:discussion-heterogeneity}

The methodological contribution of this paper is the use of cross-sectoral heterogeneity as a direct mechanism-identifying strategy, supplementing the conventional partial-identification sensitivity bounds. The conventional approach (Oster bounds, Rambachan-Roth honest-DiD) establishes that the treatment effect is robust to unobserved-confounder selection of plausible magnitude. These bounds are essential but insufficient: they confirm that an effect exists but not why it operates as it does.

The heterogeneity-test approach identifies the operating channel by exploiting cross-sectional variation in the channel-relevant parameter. The empirical-test prediction is sharper than the average-effect prediction: it requires not only that the effect exist, but that it concentrate in the cross-section where the channel-relevant parameter is highest. The sharper prediction provides identification at the L4 structural level that the L3.a + L3.b sensitivity-bound approach cannot.

The lesson generalises beyond the offtake-commitment context. Empirical-policy-evaluation studies that report only average treatment effects with partial-identification sensitivity bounds are leaving methodological value on the table: they cannot identify the operating channel separately from the existence of the effect. Heterogeneity tests against pre-registered channel-relevant cross-sectional patterns provide the additional identification.

\subsection{Policy implication: EU IF reform with offtake-commitment eligibility}
\label{sec:discussion-policy}

The policy implication of the empirical findings in this paper is concrete: capex-grant programmes that currently address the financing-constraint friction alone should incorporate demonstrable pre-FID offtake-commitment eligibility requirements. The specific reform proposal for the EU Innovation Fund is that grant eligibility be conditioned on: (i) a legally-binding offtake commitment of at least 50\% of nameplate capacity, (ii) at a minimum tenor of seven years from project commissioning, (iii) with a counterparty whose creditworthiness meets specified rating thresholds, and (iv) at indexed pricing terms that pass through input-cost variation to the offtaker rather than retaining it on the producer.

The empirical case for this reform rests on the joint identification of the offtake-commitment effect magnitude documented in this paper and the EU IF null finding documented in companion work \cite{saakstra2026paper2}. The unconditioned EU IF produces a precise null because the financing-constraint friction is not binding in the contemporary capital-abundant environment; the offtake-commitment-conditioned EU IF would jointly address F4 (the IF's current target) and F7 (the binding friction this paper identifies), restoring the IF to substantive policy effectiveness.

\subsection{Limitations}
\label{sec:discussion-limitations}

Three limitations warrant explicit acknowledgement. First, the offtake-commitment status is observed at the project level but not always with the contractual-detail precision that would be ideal for a fully identified analysis: in approximately 8\% of cases, we rely on trade-press confirmation rather than primary-document disclosure. Second, the sectoral classification by revenue volatility is constructed from observed long-term price variation in the relevant end-use market, but this construction is itself an empirical choice that admits alternative specifications; we report sensitivity to alternative classifications in Section~\ref{sec:results-robustness}, but a fully model-derived classification would require a structural-estimation extension. Third, the offtake-commitment status is observed at one point in time and treated as exogenous in the matching framework; a fully dynamic analysis would model the joint determination of offtake-commitment status and project-continuation status, which would require additional identifying restrictions beyond the scope of the present paper.

\subsection{Methodological positioning within the observational-causal-inference tradition}
\label{sec:p3-method-positioning}

The five-estimator triangulation strategy adopted in this paper --- LPM with rich controls, propensity-score matching, IPWRA, sector-stratified estimation, and Oster sensitivity bounds --- locates the empirical strategy within the broader Imbens-tradition framework of observational causal inference \cite{imbens2018causal} as a disciplined multi-strategy risk-management problem. The $\sigma$-channel interpretation of the offtake-commitment effect derives its theoretical motivation from the real-options treatment of irreversible investment under stochastic payoff variation; the recent extension of this machinery to carbon-pricing under multiple stochastic sources by \citet{olijslagers2021carbon} provides the contemporary theoretical anchor for the framing adopted here. The reported channel-attribution discussion, in which the $\sigma$-channel is identified as the dominant operating mechanism with explicit acknowledgement that the channel cannot be empirically structurally separated from the $\mu$- and $\eta$-channel alternatives in the available observational sample, is congruent with the methodological programme of \citet{blasques2024mitigating} on \emph{mitigating estimation risk} through convergent identification strategies with explicit limits.

\section{Conclusion}
\label{sec:conclusion}

This paper has documented the offtake-commitment effect on clean-hydrogen-project cancellation probability with five convergent matching estimators (point estimate $-11.3$ to $-13.2$ percentage points), exceptional robustness to unobserved selection (Oster $\delta_{\text{null}} = 20.23$), and direct mechanism identification via cross-sectoral heterogeneity (high-$\sigma$ block effect approximately ten times the low-$\sigma$ block effect, consistent with the $\sigma$-channel prediction of the real-options framework). The combined L3.a + L3.b + L4 identification is the strongest in the dissertation sample.

The methodological contribution is the use of pre-registered cross-sectoral heterogeneity tests as a direct mechanism-identifying strategy, supplementing the conventional partial-identification sensitivity-bound approach. The policy implication is concrete: capex-grant programmes such as the EU Innovation Fund should incorporate demonstrable pre-FID offtake-commitment eligibility requirements, jointly addressing the financing-constraint and counterparty-risk frictions.

Future work should extend the analysis to alternative clean-tech sectors where similar offtake-mechanism patterns may operate (battery-grade materials, direct-air-capture, sustainable aviation fuel), to dynamic models that jointly determine offtake-commitment status and project-continuation status, and to structural-estimation extensions that recover the channel-specific parameters of the underlying real-options model from project-level investment-decision data.

\bibliographystyle{plainnat}
\bibliography{paper3_references}

\end{document}
